\documentclass[aspectratio=169,11pt]{beamer}

% ============================================
% PACKAGES
% ============================================
\usepackage[utf8]{inputenc}
\usepackage[T1]{fontenc}
\usepackage[ngerman]{babel}
\usepackage{lmodern}
\usepackage{tcolorbox}
\tcbuselibrary{skins, hooks}
\usepackage{listings}
\usepackage{xcolor}
\usepackage{fontawesome5}

% Modern Font
\IfFileExists{FiraSans.sty}{\usepackage[sfdefault]{FiraSans}}{}
\renewcommand{\familydefault}{\sfdefault}

% ============================================
% COLORS & DESIGN SYSTEM
% ============================================
\definecolor{BgColor}{HTML}{FAFAFA}
\definecolor{Primary}{HTML}{2B3A42}
\definecolor{Accent}{HTML}{E84A5F}
\definecolor{Secondary}{HTML}{3F5765}
\definecolor{CodeBg}{HTML}{F0F0F0}
\definecolor{Success}{HTML}{28A745}

\setbeamercolor{background canvas}{bg=BgColor}
\setbeamercolor{title}{fg=Primary}
\setbeamercolor{frametitle}{bg=Primary, fg=white}
\setbeamercolor{structure}{fg=Accent}
\setbeamercolor{normal text}{fg=Primary}
\setbeamercolor{itemize item}{fg=Accent}

\setbeamertemplate{navigation symbols}{}
\setbeamertemplate{footline}{%
  \leavevmode%
  \hbox{%
  \begin{beamercolorbox}[wd=.333333\paperwidth,ht=2.25ex,dp=1ex,center]{author in head/foot}%
    \usebeamerfont{author in head/foot}\tiny Falko Himmel
  \end{beamercolorbox}%
  \begin{beamercolorbox}[wd=.333333\paperwidth,ht=2.25ex,dp=1ex,center]{title in head/foot}%
    \usebeamerfont{title in head/foot}\tiny Grundlagen der Praktischen Informatik
  \end{beamercolorbox}%
  \begin{beamercolorbox}[wd=.333333\paperwidth,ht=2.25ex,dp=1ex,right]{date in head/foot}%
    \usebeamerfont{date in head/foot}\insertframenumber{} / \inserttotalframenumber\hspace*{2ex} 
  \end{beamercolorbox}}%
  \vskip0pt%
}

% ============================================
% CUSTOM BOXES
% ============================================
\tcbset{
    modernbox/.style={
        enhanced,
        boxrule=0pt,
        frame hidden,
        colback=white,
        coltitle=white,
        colbacktitle=Secondary,
        fonttitle=\bfseries,
        drop lifted shadow,
        arc=4pt,
        toptitle=2mm,
        bottomtitle=2mm,
        left=2mm,
        right=2mm
    },
    alertbox/.style={
        modernbox,
        colbacktitle=Accent
    }
}

% ============================================
% CODE LISTINGS
% ============================================
\lstdefinelanguage{Haskell}{
  keywords={class, data, deriving, do, else, if, in, import, let, module, of, then, type, where, case, error},
  keywordstyle=\color{Accent}\bfseries,
  identifierstyle=\color{Primary},
  sensitive=true,
  comment=[l]{--},
  commentstyle=\color{gray}\ttfamily\itshape,
  stringstyle=\color{teal}\ttfamily,
  morestring=[b]"
}

\lstset{
    language=Haskell,
    backgroundcolor=\color{CodeBg},
    basicstyle=\ttfamily\scriptsize,
    breaklines=true,
    frame=single,
    rulecolor=\color{CodeBg},
    frameround=tttt,
    numbers=left,
    numberstyle=\tiny\color{gray},
    xleftmargin=15pt,
    tabsize=4,
    showstringspaces=false
}

% ============================================
% TITLE PAGE
% ============================================
\title{\textbf{Tutorium 6}}
\subtitle{Higher Order Functions: Map, Fold \& Filter}
\institute{Grundlagen der Praktischen Informatik}
\author{}
\date{\today}

\begin{document}

\begin{frame}[plain]
    \titlepage
\end{frame}

\begin{frame}{Inhaltsverzeichnis}
    \tableofcontents
\end{frame}

% ============================================
% THEORIE
% ============================================
\section{Lambda-Funktionen}

\begin{frame}[fragile]{Anonyme Funktionen (Lambdas)}
    Lambda-Funktionen sind anonyme Funktionen, die man ``on the fly'' erstellen kann, ohne ihnen einen Namen zu geben.
    
    \begin{tcolorbox}[modernbox, title=Syntax einer Lambda-Funktion]
\begin{lstlisting}
-- Normal:
addOne' :: Int -> Int
addOne' x = x + 1

-- Mit Lambda:
addOne :: Int -> Int
addOne x = (\x -> x + 1) x
\end{lstlisting}
    \end{tcolorbox}
    \small Startet mit \texttt{\textbackslash} (Backslash), gefolgt von Parametern, einem Pfeil \texttt{->} und dem Ausdruck.
\end{frame}

\section{Map \& Filter}

\begin{frame}[fragile]{Die \texttt{map} Funktion}
    \texttt{map} wendet eine Funktion auf \textbf{jedes} Element einer Liste an. Die Länge und Reihenfolge der Liste bleibt erhalten!
    
    \begin{tcolorbox}[modernbox, title=Beispiel: Map]
\begin{lstlisting}
-- Alle Elemente quadrieren
squareListMap :: Num a => [a] -> [a]
squareListMap xs = map (\x -> x^2) xs

-- Nur gerade Zahlen verdoppeln
doubleEvens :: [Int] -> [Int]
doubleEvens xs = map (\x -> if even x then x * 2 else x) xs
\end{lstlisting}
    \end{tcolorbox}
\end{frame}

\begin{frame}[fragile]{Die \texttt{filter} Funktion}
    \texttt{filter} prüft jedes Element gegen eine Bedingung (Prädikat) und behält nur jene, für die \texttt{True} zurückgegeben wird.
    
    \begin{tcolorbox}[modernbox, title=Beispiel: Filter]
\begin{lstlisting}
filterEvens :: [Int] -> [Int]
filterEvens xs = filter (\x -> even x) xs
-- Aufruf: filterEvens [1,2,3,4,5,6] -> [2,4,6]

filterGreaterThanFive :: [Int] -> [Int]
filterGreaterThanFive xs = filter (\x -> x > 5) xs
\end{lstlisting}
    \end{tcolorbox}
\end{frame}

\section{Fold}

\begin{frame}[fragile]{\texttt{foldr} und \texttt{foldl}}
    \texttt{fold} reduziert eine Liste auf einen einzigen Wert (z.\,B. Summe).
    \begin{itemize}
        \item \texttt{foldr} (Right): Faltet von rechts nach links. Baut den Term rekursiv auf.
        \item \texttt{foldl} (Left): Faltet von links nach rechts.
    \end{itemize}
    
    \begin{tcolorbox}[alertbox, title=Wichtig bei nicht-kommutativen Operationen (z.B. Division)]
\begin{lstlisting}
foldDivr xs = foldr (/) 1 xs
-- [2,4,8] => 2 / (4 / (8 / 1)) = 4.0

foldDivl xs = foldl (/) 1 xs
-- [2,4,8] => ((1 / 2) / 4) / 8 = 0.015625 (1/64)
\end{lstlisting}
    \end{tcolorbox}
\end{frame}

% ============================================
% LIVE DEMO
% ============================================
\section{Praxis}

\begin{frame}[plain]
    \begin{center}
        \Huge \textbf{Live Demo} \\
        \vspace{0.5cm}
        \large \faLaptopCode~ \texttt{01\_map.hs}, \texttt{02\_fold.hs}, \texttt{03\_filter.hs} \\
        \vspace{0.3cm}
        \normalsize Wir programmieren Higher Order Functions live im Code.
    \end{center}
\end{frame}

\end{document}
